%%%%%
%
\section{Interpolation lin\'eaire d'angle}

\itemstyle{Op\'erateur~:} \textsf{Inter\_Angl}
\index{Inter\_Angl}

%%%%%OBJET%%%%
\itemstyle{Objet~:}
Sur-�chantillonnage par interpolation lin\'eaire d'une variable 
angulaire exprim�e en radians.

%%%%%ARGUMENTS%%%%
\bigskip
%%%%%ENTREES%%%%
\begin{tabarg}
\multicolumn{3}{l}{\titretab {Entr\'ees}}\\
\hline
$X$ & Vecteur, {\tt float} & Param\`etre angulaire ($radian$)\\
$R_f$ & Coeff., {\tt Int} &  Rapport des fr\'equences  ($F_{sortie}\over
F_{entree}$ )\\
$Decal$ & Coeff., {\tt Int} & D\'ecalage (\textit{nombre de points}) \\
\end{tabarg}

\medskip
%%%%%SORTIES%%%%
\begin{tabarg}
\multicolumn{3}{l}{\titretab {Sorties}}\\
\hline
$X_I$ & Vecteur, {\tt float} & Angle interpol\'e($radian$)\\
\end{tabarg}


%%%%%FORMULE%%%%
\itemstyle{Formule}

%%Inter_Angl

Pour chaque $X[i]$ (\'echantillon de $X$), 
on  d�finit $a$ et $b$  : 
$a=X[i]$ et $b=X[i+1]\pm 2\pi$, tel que
$|b-a| < \pi$.

puis on calcule $R_f$ valeurs
par interpolation lin\'eaire entre $a$ et $b$.



%%%%%COMMENTAIRE%%%%
\itemstyle{Commentaire~:}

$R_f$ doit �tre sup�rieur ou �gal � 1, et le nombre d'\'echantillons
du vecteur de sortie est le produit du
nombre d'\'echantillons du vecteur d'entr\'ee par $R_f$.

L'algorithme choisi fonctionne quel que soit l'intervalle de d\'efinition
de la variable d'entr\'ee ($[-\pi,\pi]$ ou $[0,2\pi]$, \ldots), mais 
peut produire des valeurs sortant de cet intervalle.

%%%%%AUTEUR%%%%
\itemstyle{Auteur~:}
CNRM/GMEI/TRAMM

%%%%%REFERENCE%%%%
\itemstyle{R\'ef\'erences~:}
